\documentclass[11pt,a4paper]{article}
\usepackage[italian]{babel}
\usepackage[margin=2.5cm]{geometry}
\usepackage{parskip}
\usepackage{amsmath, amssymb, amsthm}
\usepackage{booktabs}
\usepackage{array}
\usepackage{xcolor}
\usepackage{needspace}
\usepackage{enumitem}
\usepackage{listings}
\usepackage{tikz}
\usepackage[most]{tcolorbox}
\usepackage[hidelinks]{hyperref}
\usetikzlibrary{decorations.pathreplacing, shapes.geometric, arrows.meta, positioning}

% Niente vedove/orfane: i paragrafi non lasciano righe isolate a fine/inizio pagina
\widowpenalty=10000
\clubpenalty=10000

% Elenchi più compatti
\setlist{itemsep=2pt, parsep=0pt, topsep=4pt, partopsep=0pt}

% --- Riquadri con bordo nero, senza numero (si spezzano tra due pagine) ---
% Uso: \begin{definizione} ... \end{definizione}
%      \begin{teorema}[Nome] ... \end{teorema}  (il nome è facoltativo)
\tcbset{nota/.style={enhanced jigsaw, breakable, colback=white,
  colframe=black, boxrule=0.5pt, sharp corners}}
\NewTColorBox{definizione}{}{nota,
  before upper={\textbf{Definizione.}\ }}
\NewTColorBox{teorema}{o}{nota,
  before upper={\textbf{Teorema\IfValueT{#1}{ (#1)}.}\ }}
\NewTColorBox{proposizione}{o}{nota,
  before upper={\textbf{Proposizione\IfValueT{#1}{ (#1)}.}\ }}
\NewTColorBox{proprieta}{o}{nota,
  before upper={\textbf{Proprietà\IfValueT{#1}{ (#1)}.}\ }}
\NewTColorBox{assioma}{o}{nota, boxrule=1pt,
  before upper={\textbf{Assioma\IfValueT{#1}{ (#1)}.}\ }}

% --- Ambienti semplici (senza riquadro) ---
% Il titolo sta in cima al blocco, anche se il contenuto inizia con un riquadro o
% con codice; e non resta mai da solo in fondo a una pagina.
% Uso: \begin{esempio}[Titolo facoltativo] ... \end{esempio}
\newcommand{\newrunin}[2]{%
  \NewDocumentEnvironment{#1}{o}{%
    \par\Needspace{4\baselineskip}\addvspace{\medskipamount}\noindent
    \textbf{#2}\IfValueT{##1}{ (##1)}\textbf{.}\ \ignorespaces}%
  {\par\addvspace{\medskipamount}}}
\newrunin{esempio}{Esempio}
\newrunin{esercizio}{Esercizio}
\newrunin{osservazione}{Osservazione}

% V e F nelle tavole di verità
\newcommand{\V}{\textsf{V}}
\newcommand{\F}{\textsf{F}}

% --- Codice C ---
% Uso: \begin{codice} ... \end{codice}      codice C numerato
%      \begin{comando} ... \end{comando}    comandi da terminale
%      \begin{risultato} ... \end{risultato} output di un programma
%      \lstinline|printf("ciao");|           codice dentro al testo
\lstdefinestyle{c}{
  language=C,
  basicstyle=\ttfamily\small,
  keywordstyle=\color{blue}\bfseries,
  commentstyle=\color{gray}\itshape,
  stringstyle=\color{red!70!black},
  numbers=left,
  numberstyle=\tiny\color{gray},
  numbersep=8pt,
  columns=flexible,
  keepspaces=true,
  breaklines=true,
  showstringspaces=false,
  tabsize=4,
  morekeywords={include, define},
  % lettere accentate nei commenti
  literate={à}{{\`a}}1 {è}{{\`e}}1 {é}{{\'e}}1 {ì}{{\`i}}1
           {ò}{{\`o}}1 {ù}{{\`u}}1 {È}{{\`E}}1 {À}{{\`A}}1
}
\lstset{style=c}
\newtcblisting{codice}{listing only, nota, left=6mm,
  listing options={style=c}}
\newtcblisting{comando}{listing only, nota,
  listing options={basicstyle=\ttfamily\small, numbers=none, language={},
    columns=flexible, keepspaces=true, breaklines=true}}
\newtcblisting{risultato}{listing only, enhanced jigsaw, breakable,
  colback=black!5, colframe=black, boxrule=0.5pt, sharp corners,
  listing options={basicstyle=\ttfamily\small, numbers=none, language={},
    columns=flexible, keepspaces=true, breaklines=true}}

% --- Diagrammi di flusso (simboli standard) ---
\tikzset{
  terminale/.style = {ellipse, draw, minimum width=2.5cm, minimum height=0.9cm},
  io/.style        = {trapezium, trapezium left angle=70, trapezium right angle=110,
                      draw, minimum width=2.5cm, minimum height=0.9cm},
  processo/.style  = {rectangle, draw, minimum width=2.5cm, minimum height=0.9cm},
  decisione/.style = {diamond, draw, aspect=2, minimum width=2.5cm},
  freccia/.style   = {thick, -{Stealth}}
}

\title{Forma indeterminata $\infty-\infty$ e teorema del confronto}
\author{Marco Cerbella}
\date{Lezione 13 -- 5 novembre 2026}

\begin{document}
\maketitle

\section{Forma indeterminata $\infty - \infty$}

Consideriamo le funzioni
\[
f_1(x) = x + x^2, \qquad f_2(x) = x + 1, \qquad f_3(x) = x + \frac{1}{x}, \qquad g(x) = x,
\]
con $\displaystyle\lim_{x\to+\infty} g(x) = +\infty$. Si ha
\[
\lim_{x\to+\infty} f_1(x) = +\infty, \qquad \lim_{x\to+\infty} f_2(x) = +\infty, \qquad \lim_{x\to+\infty} f_3(x) = +\infty \quad\left(\tfrac{1}{x} \to 0\right).
\]
Quindi $\displaystyle\lim_{x\to+\infty}\big(f_i(x) - g(x)\big)$ è una forma indeterminata $\infty - \infty$. Eppure i tre limiti hanno risultati diversi:
\begin{align*}
  \lim_{x\to+\infty}\big(f_1(x) - g(x)\big) &= \lim_{x\to+\infty}(x + x^2 - x) = \lim_{x\to+\infty} x^2 = +\infty,\\
  \lim_{x\to+\infty}\big(f_2(x) - g(x)\big) &= 1,\\
  \lim_{x\to+\infty}\big(f_3(x) - g(x)\big) &= 0 .
\end{align*}
Per questo $\infty - \infty$ è una \emph{forma indeterminata}: il risultato dipende dalle funzioni coinvolte.

\begin{esempio}[Raccoglimento]
\[
\lim_{x\to-\infty}\left(x + x^2\right) \quad (\text{F.I. } \infty - \infty): \qquad
\lim_{x\to-\infty} x^2\left(\frac{x}{x^2} + \frac{x^2}{x^2}\right) = \lim_{x\to-\infty} x^2\left(\frac{1}{x} + 1\right) = +\infty ,
\]
perché $x^2 \to +\infty$ e la parentesi tende a $0 + 1 = 1$.
\end{esempio}

\begin{esempio}[Razionalizzazione]
\[
\lim_{x\to+\infty}\left(\sqrt{1 + x} - \sqrt{x}\right) \quad (\text{F.I. } \infty - \infty).
\]
Si moltiplica e divide per $\sqrt{1 + x} + \sqrt{x}$ e si usa $(a - b)(a + b) = a^2 - b^2$:
\begin{align*}
  \lim_{x\to+\infty}\left(\sqrt{1+x} - \sqrt{x}\right)\cdot\frac{\sqrt{1+x} + \sqrt{x}}{\sqrt{1+x} + \sqrt{x}}
  &= \lim_{x\to+\infty}\frac{\left(\sqrt{1+x} - \sqrt{x}\right)\left(\sqrt{1+x} + \sqrt{x}\right)}{\sqrt{1+x} + \sqrt{x}}\\
  &= \lim_{x\to+\infty}\frac{1 + x - x}{\sqrt{1+x} + \sqrt{x}}
   = \lim_{x\to+\infty}\frac{1}{\sqrt{1+x} + \sqrt{x}} = 0,
\end{align*}
perché al denominatore $\sqrt{1+x} + \sqrt{x} \to +\infty + \infty = +\infty$.
\end{esempio}

\section{Teorema del confronto}

\begin{teorema}[del confronto]
Siano $f : D_f \to \mathbb{R}$, $h : D_h \to \mathbb{R}$, $g : D_g \to \mathbb{R}$. Se
\begin{enumerate}
  \item per $x \to c$, $f(x) \to \ell$ e $g(x) \to \ell$;
  \item esiste un intorno $I(c)$ di $c$ nel quale siano definite le tre funzioni (tranne al più nel punto $c$) e tale che
  \[
    f(x) \le h(x) \le g(x) \qquad \forall x \in I(c)\setminus\{c\},
  \]
\end{enumerate}
allora anche $h(x) \to \ell$ per $x \to c$.
\end{teorema}

Le due funzioni $f$ e $g$ «stringono» $h$: se $f(x) \to \ell$ e $g(x) \to \ell$ e $f(x) \le h(x) \le g(x)$, allora anche $h(x) \to \ell$.

\begin{proof}
Vogliamo dimostrare che $\forall\varepsilon > 0$ esiste $I(c)$ (un intervallo $]c - \delta, c + \delta[$) tale che $\forall x \in I(c)\setminus\{c\}$ si abbia $|h(x) - \ell| < \varepsilon$.

Fissiamo $\varepsilon > 0$. Poiché $\displaystyle\lim_{x\to c} f(x) = \ell$, esiste un intorno $I'(c)$ tale che $\forall x \in I'(c)\setminus\{c\}$, $|f(x) - \ell| < \varepsilon$, cioè
\[
  -\varepsilon < f(x) - \ell < \varepsilon \quad\Longrightarrow\quad -\varepsilon + \ell < f(x) < \varepsilon + \ell .
\]
Poiché $\displaystyle\lim_{x\to c} g(x) = \ell$, esiste un intorno $I''(c)$ tale che $\forall x \in I''(c)\setminus\{c\}$, $|g(x) - \ell| < \varepsilon$, cioè
\[
  -\varepsilon + \ell < g(x) < \varepsilon + \ell .
\]
Dall'ipotesi 2, $\forall x \in I(c)\cap I'(c)\cap I''(c)\setminus\{c\}$ si ha
\[
  -\varepsilon + \ell < f(x) \le h(x) \le g(x) < \ell + \varepsilon ,
\]
dove la prima disuguaglianza vale perché $x \in I'(c)$, la seconda e la terza perché $x \in I(c)$, la quarta perché $x \in I''(c)$. Posto $I'''(c) = I(c)\cap I'(c)\cap I''(c)$, per ogni $x \in I'''(c)\setminus\{c\}$ si ha
\[
  -\varepsilon + \ell < h(x) < \ell + \varepsilon \iff |h(x) - \ell| < \varepsilon .
\]
Quindi $h(x) \to \ell$ per $x \to c$.
\end{proof}

\end{document}
